\subsection{Framework}
We implemented a PyTorch-based framework that abstracts agents and environments, ensuring they are interchangeable and easily configurable for training.
In a configuration file, the environment\footnote{Our framework supports the OpenAI Gymnasium environment and the Hockey environment.}, the agent type, its training parameters, hyperparameters, and opponents can be specified.

The framework also enables (asynmmetric) self-play of the agents, where the trained agent is periodically replaced by one of its opponents at a specified frequency. The rule-based agents of the hockey environment can also be used as opponents but are not considered as training candidates.


% A major challenge in reinforcement learning is the formulation of a appropriate reward function for training. Thus the reward for scoring a goal in the hockey environment is sparse, other reward signals have to be considered to enable fast convergence and good exploration.
% Therefore we implemented a custom reward function, rewarding the aspects: scoring or getting goals, the closeness to the puck, the puck direction, touching the puck, blocking the puck and staying in the goal. We computed the total reward by summing up the rewards of the different aspects, each weighted by a factor that can be configurated. For the evaluation step just the reward for scoring a goal was used.

The leaderboard of our framework ranks agents dynamically based on their performance using the Placket-Luce rating system \cite{af5079a1-8ca5-3727-a405-0a82390327b7}. There, each agent has a skill rating average $\mu$ and variance $\sigma$, which are updated after every match. Games are simulated, and winners gain rating points while losers may lose them. Agents are paired based on a probability distribution over their skill levels, ensuring fair matchups. To ensure a diverse tournament, we first sample 4 agents and assign them probabilities based on the softmax of their $\sigma$. Based on them, the final 2 agents are sampled. Agent names, types, ratings, and average scores are displayed on the leaderboard.
